\documentclass{subfiles}
\begin{document}
    Our description of move-to-front and run-length-encoding will be the following:
        we describe MTF in details, presenting the core idea and an example;
        for RLE we simply give the general idea.

    Introduced in \cite{bentleySTW86}, MTF is a simple encoding based on frequency.
        Formally, let \(S\) be a text and let \(A\) be the alphabet on which its defined.
        It follows, that \(S\)'s characters satisfy some lexicographic order.
        Using this fact, MTF encodes the \(i\)-th symbol in \(S\) by the 
        number of distinct symbols preceeding it in \(A\), 
        then moving said character to the front of \(A\).
        The process is repeated until \(S\) is scanned completely.

    The core idea is that, if \(S\) is localy homogeneous, as in the case of BWT,
        then the output of the encoding will be a sequence of small integers,
        which can be encoded in turns by some universal encoder of integers.

    Decoding is easy as well; since the alphabet is known by both parties,
        we proceed as done in encoding. Precisely, we decode the symbol as 
        the \((i + 1)\)-th symbol in \(A\) and move that character to the front.

    Perfomance-wise it can be proved that MTF behaves slightly worse then Huffman 
        in the worst case, and a lot better in the best one.

    \begin{example*}
        Consider the string \(w = smmihtt\textdollar ecaa = bwt(mathematics\textdollar)\). 
            The alphabet is \(A = \set{\textdollar, a, c, e, h, i, m, s, t}\) and the lexicographic 
            order is the usual one \((\textdollar \prec a \prec \cdots \prec t)\). 

        We encode \(s\) as \(7\), and move it to the fron of \(A\);
            \(m\) is then encoded as \(7\), and so on. 
            Once \(w\) has be read completely, the output of MTF is the sequence
            \((7, 7, 0, 7, 7, 8, 0, 5, 8, 8, 8, 0)\). 
            Note that \(w\) has little-to-no run, thus the ``big'' integers are expected.
    \end{example*}

    Run-length-encoding is somewhat simple, and several implementations can be found
        in the literature. The idea is that, if a character \(x\) appears \(n\)
        consecutive times, then it can be encoded as \(nx\).
        How to encode \(nx\) depends on the implementations,
        some use pairs eg. \((x, n)\); other encode \(n\) replacing the 1s in 
        its binary representation with \(x\). 
\end{document}
